\documentclass[../DoAn.tex]{subfiles}
\begin{document}

\section{Đặt vấn đề}
\label{section:1.1}

Trong bối cảnh đô thị hóa nhanh và mật độ phương tiện gia tăng, bài toán tổ
chức và quản lý giao thông ngày càng trở nên phức tạp, đòi hỏi các phương pháp
phân tích có căn cứ dữ liệu để hỗ trợ quy hoạch, đánh giá phương án tổ chức
giao thông và kiểm chứng các giả thuyết kỹ thuật. Trên thực tế, việc thử nghiệm
trực tiếp ngoài hiện trường thường tốn kém, khó kiểm soát biến số, và có thể
tiềm ẩn rủi ro về nhân mạng. Vì vậy, mô phỏng giao thông trở thành một công cụ
quan trọng, giúp tái hiện nhiều kịch bản khác nhau trong điều kiện có thể kiểm
soát, từ đó hỗ trợ quá trình ra quyết định.

Trong nhóm các công cụ mô phỏng, \gls{SUMO} (Simulation of Urban MObility)~\cite{sumo2012,sumo_web} là
một nền tảng mô phỏng vi mô được sử dụng rộng rãi, cho phép mô tả mạng đường,
luồng phương tiện và cơ chế điều khiển (như tín hiệu đèn) với mức độ chi tiết
cao. Kết quả mô phỏng có thể cung cấp dữ liệu theo thời gian về vị trí, tốc độ,
hướng di chuyển của phương tiện, cũng như các thông tin liên quan đến trạng thái
mạng. Tuy nhiên, dữ liệu đầu ra này thường có khối lượng lớn và giàu tính kỹ
thuật, khiến việc quan sát và khai thác hiệu quả trở thành một thách thức, đặc
biệt khi người dùng cần đánh giá trực quan cấu trúc không gian hoặc diễn biến ở
các khu vực có hình học phức tạp.

Việc quan sát kết quả mô phỏng bằng giao diện 2D truyền thống thường đáp ứng tốt
các thao tác cơ bản, nhưng vẫn gặp hạn chế khi cần: (i) nhận thức không gian tại
các nút giao nhiều nhánh hoặc có tổ chức làn phức tạp; (ii) theo dõi đồng thời
nhiều đối tượng trong bối cảnh mật độ cao; và (iii) trình bày kết quả cho nhóm
người dùng không chuyên hoặc trong các hoạt động trao đổi, thuyết minh phương
án. Bản chất của góc nhìn 2D là góc nhìn từ trên cao: nó cho phép người quản lý
giám sát giao thông thấy được toàn cảnh mạng đường và dòng phương tiện, nhưng
lại không tái hiện được góc nhìn của người trực tiếp tham gia giao thông --- tức
là cách mà đường sá, biển báo, tín hiệu và các phương tiện khác hiện ra trước
mắt một người đang di chuyển trong mạng. Trong các trường hợp này, góc nhìn 3D
có thể giúp người quan sát nắm bắt tốt hơn mối quan hệ không gian giữa đường,
làn, phương tiện và tín hiệu điều khiển, đồng thời tạo điều kiện cho các hình
thức tương tác và trải nghiệm theo góc nhìn người tham gia giao thông.

Mặc dù vậy, các giải pháp hiển thị 3D trong thực tế thường gặp những rào cản như
phụ thuộc vào bản đồ dựng sẵn (tĩnh), cần nhiều thao tác thủ công khi thay đổi
kịch bản, hoặc khó mở rộng khi số lượng đối tượng và loại đối tượng tăng lên.
Ngoài ra, phần lớn giải pháp chỉ dừng ở mức quan sát thụ động, chưa cho phép
người dùng tương tác trực tiếp với phương tiện trong cảnh để khảo sát các tình
huống giả định. Việc chuyển dữ liệu mô phỏng sang môi trường 3D cũng phát sinh
các vấn đề kỹ thuật như ánh xạ hình học từ mô tả mạng đường (đa tuyến) sang mô
hình cảnh, chuyển đổi hệ tọa độ, đồng bộ thời gian giữa bước mô phỏng và khung
hình hiển thị, và đảm bảo hiệu năng khi cập nhật số lượng lớn đối tượng theo
thời gian.

Do đó, việc xây dựng một hệ thống hiển thị giao thông 3D dựa trên dữ liệu đầu ra
của \gls{SUMO}, có khả năng tự động dựng cảnh theo bản đồ mô phỏng, cập nhật
trạng thái theo thời gian và cho phép tương tác, là cần thiết để nâng cao hiệu
quả quan sát, phân tích và truyền đạt kết quả mô phỏng. Hệ thống này đóng vai
trò cầu nối giữa dữ liệu mô phỏng (mang tính kỹ thuật) và nhu cầu khai thác trực
quan, giúp người dùng tiếp cận, kiểm chứng và trình bày kịch bản giao thông một
cách trực quan và nhất quán.

\section{Mục tiêu và phạm vi đề tài}
\label{section:1.2}

Từ các vấn đề nêu ở phần \ref{section:1.1}, đồ án hướng tới mục tiêu xây dựng một
hệ thống hiển thị mô phỏng giao thông trong không gian 3D, trong đó dữ liệu được
lấy trực tiếp từ quá trình chạy mô phỏng của \gls{SUMO}. Trọng tâm của hệ thống
là hình thành một quy trình chuyển đổi dữ liệu mạng đường và dữ liệu mô phỏng
sang các đối tượng 3D trong Unity một cách nhất quán, đồng thời bổ sung khả năng
tương tác để người dùng không chỉ quan sát mà còn có thể can thiệp vào diễn biến
giao thông.

Về mục tiêu cụ thể, đồ án tập trung vào bốn nhóm mục tiêu chính. Thứ nhất, xây
dựng cơ chế dựng mạng đường trong không gian 3D từ dữ liệu của \gls{SUMO}, bảo
đảm các đối tượng được dựng bám sát dữ liệu nguồn; mạng đường có thể được sinh
tự động từ lưới mê cung phục vụ kiểm thử hoặc nhập từ bản đồ thực tế
OpenStreetMap~\cite{osm}. Thứ hai, xây dựng cơ chế thu thập dữ liệu động theo thời gian
thông qua \gls{TraCI}~\cite{traci2008}, nhằm lấy trạng thái của phương tiện, đèn tín hiệu và
người đi bộ theo từng bước mô phỏng. Thứ ba, thiết kế cơ chế cập nhật và hiển
thị trạng thái trong Unity theo thời gian, bảo đảm chuyển động thể hiện đúng quỹ
đạo tổng quát và duy trì hiệu năng ổn định khi số lượng đối tượng tăng. Thứ tư,
xây dựng khả năng tương tác cho phép người dùng chiếm quyền điều khiển một
phương tiện và lái thủ công bằng vật lý, xử lý va chạm và trả quyền điều khiển
để phương tiện hòa trở lại dòng giao thông do \gls{SUMO} quản lý.

Trong phạm vi của đề tài, hệ thống tập trung vào các thành phần cốt lõi gồm: (i)
dựng cảnh 3D từ dữ liệu mạng đường, bao gồm các đoạn đường, các yếu tố hình
học phục vụ quan sát, và khối công trình (toà nhà) khi nhập từ bản đồ
OpenStreetMap; (ii) biểu diễn phương tiện và người đi bộ như các thực thể
di chuyển theo trạng thái thời gian; (iii) biểu diễn các tín hiệu điều khiển ở
mức phù hợp với dữ liệu có thể truy xuất; (iv) chuẩn hóa, ánh xạ dữ liệu sang hệ
tọa độ của Unity để bảo đảm tính nhất quán; và (v) hỗ trợ tương tác lái xe ở mức
một người dùng cho mỗi phiên mô phỏng. Phần đánh giá được thực hiện thông qua
các kịch bản mô phỏng mẫu nhằm kiểm chứng tính đúng đắn của việc ánh xạ dữ liệu,
khả năng tái hiện diễn biến theo thời gian và tính ổn định của hiển thị.

Ngoài phạm vi của đồ án là các bài toán tối ưu điều khiển giao thông, hiệu
chỉnh hay tham số hóa mô hình mô phỏng của \gls{SUMO}, cũng như hỗ trợ nhiều
người dùng đồng thời trong cùng một phiên. Trong phạm vi này, hệ thống đóng vai
trò là tầng trực quan hóa, tương tác và hỗ trợ phân tích dựa trên dữ liệu mô
phỏng sẵn có, qua đó tăng khả năng quan sát và trình bày kịch bản thay vì can
thiệp vào thuật toán mô phỏng.

\section{Định hướng giải pháp}
\label{section:1.3}

Để hiện thực hóa mục tiêu ở phần \ref{section:1.2}, đồ án lựa chọn hướng tiếp cận
dựa trên chuỗi công cụ: \gls{SUMO} dùng để mô phỏng giao thông vi mô và sinh dữ
liệu; Python kết hợp \gls{TraCI} dùng để truy xuất trạng thái mô phỏng theo thời
gian; Unity dùng để dựng cảnh, hiển thị 3D và xử lý tương tác. Lựa chọn này phù
hợp vì \gls{SUMO} là phần mềm mã nguồn mở, được sử dụng rộng rãi trong giới
nghiên cứu và học thuật, đồng thời cung cấp dữ liệu mô phỏng có cấu trúc rõ ràng
và phổ biến; trong khi Unity là môi trường phát triển 3D mạnh, hỗ trợ dựng cảnh,
quản lý đối tượng, camera, vật lý và tương tác, qua đó nâng cao khả năng quan sát
và can thiệp trực quan. Sự kết hợp giữa hai nền tảng này hướng tới phục vụ đồng
thời hai nhóm người dùng: nhóm nghiên cứu và học tập vốn quen thuộc với \gls{SUMO}
nay có thêm khả năng quan sát, trình bày kịch bản một cách trực quan; và nhóm
phát triển trò chơi (game) vốn quen thuộc với Unity nay có thể tận dụng dữ liệu
giao thông được mô phỏng sát thực tế làm nền tảng cho các sản phẩm tương tác.

Trên cơ sở đó, giải pháp được triển khai theo một quy trình tổng quát gồm các
thành phần chính. Thứ nhất, module chuẩn bị dữ liệu sinh mạng đường (từ mê cung
hoặc từ OpenStreetMap) cùng tuyến đường cho phương tiện và người đi bộ, tạo
thành kịch bản mô phỏng đầu vào cho \gls{SUMO}. Thứ hai, module thu thập dữ liệu
động kết nối với phiên mô phỏng qua \gls{TraCI} để lấy trạng thái theo từng bước,
sau đó chuẩn hóa thành một định dạng trao đổi thống nhất và truyền sang Unity qua
TCP/IP. Thứ ba, module hiển thị trong Unity nhận dữ liệu, ánh xạ hệ tọa độ, tạo
và cập nhật các đối tượng 3D, đồng thời nội suy chuyển động để hiển thị mượt. Thứ
tư, module tương tác cho phép người dùng chiếm quyền điều khiển một phương tiện,
lái bằng vật lý, xử lý va chạm để sinh trạng thái ``xác xe'', và trả quyền điều
khiển để xóa xe khỏi mô phỏng ngay lập tức.

Trong quá trình triển khai, một yêu cầu quan trọng là bảo đảm tính nhất quán của
dữ liệu giữa các module: dữ liệu hình học mạng đường cần dùng chung chuẩn tọa độ
với dữ liệu phương tiện; các đối tượng trong Unity cần có vòng đời phù hợp với
trạng thái mô phỏng (tạo mới khi xuất hiện, cập nhật khi di chuyển, loại bỏ khi
rời mạng). Đồng thời, để bảo đảm khả năng mở rộng, giải pháp ưu tiên tách biệt
phần thu thập/chuẩn hóa dữ liệu (Python) khỏi phần hiển thị (Unity), giúp giảm
phụ thuộc và thuận tiện khi thay đổi kịch bản hoặc nền tảng hiển thị.

Đóng góp chính của đồ án là đề xuất và hiện thực một quy trình chuyển đổi dữ liệu
mô phỏng từ \gls{SUMO} sang môi trường hiển thị 3D trong Unity theo hướng tự
động, giảm thao tác thủ công khi thay đổi kịch bản; đồng thời bổ sung khả năng
tương tác lái xe, xử lý va chạm và trả quyền điều khiển vốn ít gặp ở các giải
pháp hiển thị thông thường. Kết quả thử nghiệm trên các kịch bản mẫu cho thấy hệ
thống có thể dựng được mạng đường từ dữ liệu đầu vào, biểu diễn chuyển động của
phương tiện và người đi bộ một cách ổn định, và cho phép người dùng can thiệp
vào diễn biến giao thông, tạo cơ sở cho việc quan sát và phân tích trực quan.

\section{Bố cục đồ án}
\label{section:1.4}

Phần còn lại của báo cáo đồ án tốt nghiệp này được tổ chức như sau. Chương 2
trình bày khảo sát hiện trạng và phân tích yêu cầu của hệ thống, trong đó so
sánh các công cụ hiển thị mô phỏng hiện có, xác định các tác nhân, mô tả tổng
quan chức năng qua biểu đồ use case và đặc tả chi tiết các use case quan trọng
cùng các yêu cầu phi chức năng. Chương 3 giới thiệu nền tảng lý thuyết và các
công nghệ được sử dụng, bao gồm \gls{SUMO}, \gls{TraCI}, Python và Unity, đồng
thời phân tích vai trò của từng thành phần và lý do lựa chọn so với các phương án
thay thế. Chương 4 trình bày phân tích thiết kế, triển khai và đánh giá hệ thống,
tập trung vào kiến trúc tổng thể, các module chính, cơ chế đồng bộ dữ liệu và
kết quả thực nghiệm trên các kịch bản mẫu. Chương 5 trình bày các giải pháp kỹ
thuật và đóng góp nổi bật của đồ án, trong đó có cơ chế giao tiếp thời gian thực,
cơ chế chiếm quyền điều khiển và xử lý va chạm, cũng như cơ chế trả quyền và lưu
trữ phiên mô phỏng. Cuối cùng, Chương 6 tổng kết các kết quả đạt được, các đóng
góp chính của đồ án và đề xuất hướng phát triển trong tương lai.

\end{document}
